% Generated draft: requires mathematical and source review.
\documentclass[11pt]{article}
\usepackage[margin=25mm]{geometry}
\usepackage{fontspec}
\setmainfont{Latin Modern Roman}
\usepackage{amsmath,amssymb,mathtools,unicode-math}
\setmathfont{Latin Modern Math}
\usepackage{enumitem,needspace,xurl,hyperref}
\setlength{\parindent}{0pt}
\setlength{\parskip}{5pt}
\setlength{\emergencystretch}{3em}
\newcommand{\N}{\mathbb{N}}
\newcommand{\Z}{\mathbb{Z}}
\newcommand{\Q}{\mathbb{Q}}
\newcommand{\R}{\mathbb{R}}
\newcommand{\C}{\mathbb{C}}
\newcommand{\F}{\mathbb{F}}
\newcommand{\A}{\mathbb{A}}
\newcommand{\PP}{\mathbb{P}}
\newcommand{\E}{\mathbb{E}}
\newcommand{\eps}{\varepsilon}
\newcommand{\dd}{\,\mathrm{d}}
\begin{document}
\section*{UnsolvedMath ID: 2306112}
\textbf{Source problem code:} catalogue entry
\subsection{Definitions and mathematical statement}
Let $A$ conventionally be the normalized analytic functions $f(z)=z+\sum_{n\geq2}a_nz^n$ on $\mathbb D$, let $S\subset A$ be the univalent class, and let $St$ be the starlike class. For two analytic series, $f*g$ denotes coefficientwise Hadamard multiplication. The source defines the dual of a normal family $\mathcal F\subset A$ by
\[
 \mathcal F^*=\{g\in A:(g*f)(z)\ne0\text{ for every }f\in\mathcal F
                              \text{ and }0<|z|<1\}.
\]
In this problem $S^*$ denotes this \emph{dual of the univalent class}, not the common alternative notation for starlike functions; the source uses $St$ for the latter.

The printed definition of the $\Sigma_\delta(f)$ neighborhood has unmatched absolute-value/parenthesis markup. Its apparently intended symmetric version, with $h=g-f$, is
\[
 \Sigma_\delta(f)=\left\{g\in A:
 \left|h'(z)-\frac{h(z)}z\right|+
 \left|h'(z)+\frac{h(z)}z\right|<2\delta
 \text{ for every }z\in\mathbb D\right\}. \tag{1}
\]
At $0$ both expressions have removable value $0$. Formula (1) is conditional on recovering the damaged second absolute value as shown; it is not an authenticated silent correction. The source also alternates $\Sigma$ and $\sum$ for the neighborhood symbol.

\textbf{Question (T. Sheil-Small and E. M. Silvia).} Is
\[
 \Sigma_1(\operatorname{id})\subset S^*,\qquad \operatorname{id}(z)=z?
\]
Under (1), this asks whether every $g\in A$ obeying the displayed symmetric inequality with $f(z)=z$ satisfies $(g*u)(z)\ne0$ for every $u\in S$ and $0<|z|<1$.

The source reports $N_\delta(f)\subset\Sigma_\delta(f)$, $\Sigma_1(\operatorname{id})\subset St$, and $\Sigma_{1/4}(f)\subset St$ for convex $f$, attributing the latter neighborhood results to Sheil-Small and Silvia. The neighborhood $N_\delta$ and additional notation are referred to Problem 6.111, which is absent. It also states that the Bieberbach conjecture is equivalent to $N_1(\operatorname{id})\subset S^*$ and recalls $S^*\subset St$, referring to de Branges and Ruscheweyh without complete citations. These contextual claims do not identify the larger $\Sigma$-neighborhood with $N$, and no conclusion about the proposed inclusion is inferred from them.

The supplied summary labels the question OPEN-TRIAGE, says it was open in the 2018 edition, and reports no definitive resolution located.
\subsection{Short English statement}
Does the indicated symmetric derivative neighborhood of the identity lie in the convolution dual of the normalized univalent class? Here $S^*$ means a dual class and $St$ means starlike. The neighborhood formula has damaged absolute-value markup, and its repaired symmetric reading needs source confirmation.
\subsection{Sources}
\begin{itemize}
\item Local UnsolvedMath registry, permanent ID 2306112, source code catalogue entry, ``Research Problems in Function Theory --- Problem 6.112''; attribution: T. Sheil-Small and E. M. Silvia. Supplied status open, classification OPEN-TRIAGE, updated 2026-08-06.
\item Hayman, \emph{Research Problems in Function Theory} (2018), Problem 6.112, with cross-reference to Problem 6.111, \url{https://arxiv.org/abs/1809.07200}. The record reports source-TeX extraction and access on 2026-07-29; the 2018 update was retained for status triage. Current status and rights require review; the record describes publicly available arXiv source TeX.
\item The source names Sheil-Small and Silvia for neighborhood inclusions and refers to Ruscheweyh for convolution duality and de Branges for the Bieberbach conjecture. Complete citations are absent.
\end{itemize}
\end{document}
